\section{From an arbitrary weak exit to countably many tests}
\label{ex:section}

We first work in a smooth, time-oriented, globally hyperbolic
four-dimensional spacetime $(M,g)$ with smooth spacelike Cauchy
surface $\Sigma$.  Suppose that $M$ is identified with its proper
open image in a smooth ambient manifold $(\widetilde M,\widetilde g)$,
where $\widetilde g$ is continuous, nondegenerate, and belongs to
$W^{1,2}_{\mathrm{loc}}$.  Thus its Christoffel array belongs to
$L^2_{\mathrm{loc}}$ in every smooth chart.  Its restriction to $M$
agrees almost everywhere with the smooth Christoffel array of $g$.
All geodesics and Hamiltonian flows below are those of the smooth
metric on the full $M$.
The compact-diamond and first-exit-graph arguments are related to
\cite[Proposition~2.6 and Lemma~2.12]{Sbierski2022}. We give their
proofs here and then average smooth geodesics in phase space to obtain
the positive-measure family required at the weaker connection regularity.

\Needspace{8\baselineskip}
\subsection{Intrinsic diamonds and a first-exit graph}

\begin{lemma}[Causal homotopies with interior endpoints]
\label{ex:homotopy}
Let $p,q\in M$ and let
$H:[0,1]\times[0,1]\longrightarrow\widetilde M$ be continuous.
Suppose that every $H(r,\cdot)$ is a locally Lipschitz future causal
path from $p$ to $q$, with causal tangent almost everywhere, and that
$H(0,[0,1])\subset M$. Then the entire image of $H$
lies in $M$.
\end{lemma}

\begin{proof}
Let $A=\{r:H(r,[0,1])\subset M\}$.  Compactness of the parameter
interval and openness of $M$ imply that $A$ is relatively open.
For $r\in A$, the path is an intrinsic causal path in $M$, so its
image lies in the single compact intrinsic diamond
$K=J_M^+(p)\cap J_M^-(q)$.  The image of $K$ in the Hausdorff
manifold $\widetilde M$ is compact and hence closed.  If
$r_n\in A$ and $r_n\to r$, continuity gives
$H(r,t)\in K$ for every $t$.  Hence $A$ is also closed, and it is
nonempty.  This proves the assertion.  Global hyperbolicity has
been used only inside $M$.
\end{proof}

\begin{lemma}[A local first-exit graph]
\label{ex:graph}
Suppose that a future timelike $C^1$ curve exits $M$ as in
Definition~\ref{intro:extension}.  There is a precompact ambient coordinate chart
with coordinates $(s,z)\in\R\times\R^3$, a constant $c>0$, a
bottom height $-d_0<0$, a spatial ball $D$, and a Lipschitz
function $f:D\to\R$, with Lipschitz constant $L<\infty$, having the
following properties:
\begin{enumerate}[label=(\roman*)]
\item $s$ is temporal, and every vector $(1,\eta)$ with
      $|\eta|<c$ is uniformly future timelike;
\item $(-d_0,z)\in I_M^+(\Sigma)$ for all $z\in D$;
\item the vertical segment $\{(s,z):-d_0\le s<f(z)\}$ lies in
      $M$, while $(f(z),z)\in\partial M$;
\item $f(0)=0$ and, for some $0<\delta<d_0/4$,
      $-\delta<f(z)\le\delta$.  The rectangle containing these
      segments has an upper face above $2\delta$.
\end{enumerate}
The metric, its inverse, the positive spatial block $(g_{ij})$,
and the timelike normalization on any closed narrower cone have
uniform bounds and uniform nondegeneracy on this rectangle.
\end{lemma}

\begin{proof}
First the given exiting curve has a late segment in
$I_M^+(\Sigma)$.  Extend it to the past by a past-inextendible
timelike curve in $M$.  Its future end has no endpoint in $M$,
since its ambient limit lies on $\partial M$.  The resulting
inextendible timelike curve meets the Cauchy surface, and its
remaining future segment lies in $I_M^+(\Sigma)$.

We can replace a late part by a straight timelike first-exit
segment.  To see this without any assertion about the ambient
causal structure, choose a small chart about the given boundary
point $b$ and a point $q$ slightly to its future on a fixed,
uniformly timelike coordinate line.  By continuity of
$\widetilde g$, for a sequence of late interior points
$x_n\to b$ all straight segments $[x_n,q]$ in this chart are
uniformly future timelike.  If one has a first exit, use it.
Otherwise every such segment, including its endpoint $q$, lies
in $M$.  For a fixed earlier point $p$ on the exiting curve,
all sufficiently late $x_n$ then lie in
$J_M^+(p)\cap J_M^-(q)$.  Compactness of this intrinsic diamond
would put their ambient limit $b$ in $M$, a contradiction.
The starting point of the chosen first-exit segment is in
$I_M^+(\Sigma)$.

Translate its first exit to $(0,0)$ and make its direction
vertical by a linear coordinate change.  Choose the other
coordinate vectors in the spacelike orthogonal hyperplane of
that direction at the exit.  After shrinking the chart,
$\partial_s$ is timelike, $\dd s$ is temporal, and a fixed cone
$|\eta|<c$ has the asserted strict timelike margin.  The closure
of the chart is taken inside a larger smooth coordinate chart;
continuity and nondegeneracy give all the claimed uniform
algebraic metric bounds.

Choose $d_0>0$ so that the central segment
$[-d_0,0)\times\{0\}$ is the final part of the selected straight
segment.  Its bottom has a neighborhood in $I_M^+(\Sigma)$.
Take $0<\delta<d_0/4$.  For $z$ in a sufficiently small ball,
define the first-exit height by
\[
 f(z)=\inf\{s\in[-d_0,2\delta]:(s,z)\notin M\},
 \qquad \inf\varnothing=2\delta.
\]
Compact openness along $[-d_0,-\delta]\times\{0\}$ permits
shrinking the ball so that $f(z)>-\delta$.  Write its radius as
$R$ and shrink it further so that $R/\delta<c/4$.
If $f(z)>\delta$, both $p_z=(-d_0,z)$ and
$q_z=(\delta,z)$, and their entire vertical joining segment,
are in $M$.  The broken path
\[
                p_z\longrightarrow (0,0)\longrightarrow q_z
\]
has spatial slopes less than $c/2$.  Interpolate its spatial
graph linearly with that of the vertical segment, using $s$ as
parameter.  This is a fixed-endpoint homotopy of future timelike
paths in the chart.  \Cref{ex:homotopy} forbids the member
through $(0,0)\notin M$.  Thus $f(z)\le\delta$; in particular
the clipping at $2\delta$ is inactive.  Its definition now
implies $(f(z),z)\in\partial M$, even if a vertical line
reenters $M$ at a later height.

Finally shrink $R$ so that $2R/(d_0-\delta)<c/2$, and fix
$L>4/c$.  If $f(z')>f(z)+L|z-z'|$, select a height $b'$ with
\[
       f(z)+2|z-z'|/c<b'<f(z').
\]
The path from $(-d_0,z')$ through $(f(z),z)$ to $(b',z')$
has first-leg slope less than $2R/(d_0-\delta)<c/2$ and
second-leg slope less than $c/2$.  Its endpoints and the
reference vertical segment lie strictly below the first exit
on the higher vertical.  Linear interpolation of spatial
graphs again contradicts \Cref{ex:homotopy}.  Interchanging
$z,z'$ proves $|f(z)-f(z')|\le L|z-z'|$.  Piecewise smooth
timelike paths suffice in these homotopies; the strict slope
inequalities also permit smoothing the corners.  The common
bottom height is what keeps the first leg uniformly inside
the cone.
\end{proof}

\subsection{Averaging near the graph}

\begin{lemma}[Boxes with small normalized connection mass]
\label{ex:boxes}
Fix $K>1$ and a ball $D'\Subset D$ from \Cref{ex:graph}.
There are $l_n\downarrow0$ and $z_n\in D'$ such that the boxes
\[
 Q_l(z)=\{(s,y):|y-z|<Kl,\ |s-f(z)|<Kl\}
\]
lie in one fixed precompact subchart and satisfy
\begin{equation}\label{ex:small-box}
        l_n^{-3}\int_{Q_{l_n}(z_n)}
                   |\widetilde\Gamma|\,\dd s\,\dd y
           \longrightarrow0.
\end{equation}
\end{lemma}

\begin{proof}
On the compact chart region, $L^2$ integrability implies
$L^1$ integrability.  All the boxes with $z\in D'$ and small
$l$ lie in a fixed compact region, and their union is contained
in a layer
\[
       \mathcal L_l=\{(s,y):y\in D'',\quad
                    |s-f(y)|\le K(1+L)l\},
        \qquad D'\Subset D''\Subset D.
\]
For each $(s,y)$, the set of possible centers $z$ has volume
at most $C_Kl^3$.  Tonelli's theorem therefore gives
\[
 \int_{D'} l^{-3}\int_{Q_l(z)}|\widetilde\Gamma|
                   \,\dd s\,\dd y\,\dd z
       \le C_K\int_{\mathcal L_l}|\widetilde\Gamma|.
\]
The graph of $f$ has four-dimensional Lebesgue measure zero,
by Fubini.  Absolute continuity of the $L^1$ integral makes
the right side tend to zero.  Choosing a center with no more
than the average normalized mass proves the assertion along
any sequence of sufficiently small scales.
\end{proof}

\begin{proposition}[A positive family supplied by every weak exit]
\label{prop:exit-family}
Under the hypotheses of \Cref{ex:graph}, there is a measurable
set $S\subset T\Sigma$ of positive smooth section measure such
that every future unit geodesic with seed in $S$ has finite
maximal future proper lifetime in the full $M$ and an ambient
endpoint on $\partial M$.  The endpoints lie in one compact
subset of an ambient chart.  Moreover, there are finitely many
precompact ambient charts $\psi_j:U_j\to V_j\subset\R^4$ and
constants $C<\infty$, $m>0$ with the following properties:
\begin{enumerate}[label=(\roman*)]
\item on each entire $V_j$, the metric and inverse metric have
      component norms at most $C$, and
      $\|\widetilde\Gamma_j\|_{L^2(V_j)}\le C$;
\item for every seed in $S$ and every proper time from its
      intersection with $\Sigma$ to its full future endpoint,
      excluding only the endpoint itself, at least one $j$
      satisfies
      \begin{equation}\label{ex:chart-witness}
      \operatorname{dist}\bigl(\psi_j(\gamma(\tau)),
                                  \R^4\setminus V_j\bigr)\ge m,
      \qquad
      |(\psi_j)_*\dot\gamma(\tau)|\le C.
      \end{equation}
\end{enumerate}
Here the section measure is the pullback of canonical volume
on $T^*\Sigma$ under $w\mapsto h(w,\cdot)$, where $h$ is the
induced metric.  In particular it has the same null sets as
every smooth positive measure on $T\Sigma$.
\end{proposition}

\begin{proof}
Use the graph and chart of \Cref{ex:graph}.  Constants in this
proof depend on this fixed chart, but not on the small scale
$l$.  Choose small numbers $\rho_0,\varepsilon_0>0$ so that
\begin{equation}\label{ex:cone-choice}
 L\rho_0<\tfrac14,
 \qquad 2\varepsilon_0+C_\beta(\rho_0+2\varepsilon_0)<c,
\end{equation}
where $C_\beta$ bounds the derivative of a fixed smooth bump
on $(-1,1)$ that equals one at zero.  We may also require
$\rho_0+4\varepsilon_0<1$.  Fix $K>3$ in \Cref{ex:boxes}.
For one of its boxes, write $z_c$ for the center and set
\[
 s_c=f(z_c),\qquad s_0=s_c-l,\qquad s_1=s_c+l.
\]
Every point $(s_0,z)$ with $|z-z_c|<\rho_0l$ is strictly
below $f(z)$ and above the bottom, for small $l$.  It is in
$I_M^+(\Sigma)$, since its vertical segment from the bottom
lies in $M$ and is future timelike.

\emph{The initial phase volume.}
At such a point parameterize initial unit velocities by
slopes $\eta\in\R^3$, $|\eta|<\varepsilon_0$:
\[
  U=N(1,\eta),\qquad
  N=\bigl[-\widetilde g((1,\eta),(1,\eta))\bigr]^{-1/2},
  \qquad
  \pi_i=N(\widetilde g_{i0}+\widetilde g_{ij}\eta^j).
\]
On $|\eta|\le2\varepsilon_0$, $N$ has positive uniform lower
and upper bounds.  Differentiation in the slope variables,
which uses no derivative of the ambient metric, gives
\begin{equation}\label{ex:momentum-jacobian}
 \frac{\partial\pi_i}{\partial\eta^j}
 =N\widetilde g_{ij}
  +N^3(\widetilde g_{i0}+\widetilde g_{ik}\eta^k)
       (\widetilde g_{j0}+\widetilde g_{jk}\eta^k).
\end{equation}
The spatial block is uniformly positive definite.  Thus this
Jacobian is uniformly positive definite and has determinant
bounded above and away from zero.  Its positivity also makes
$\eta\mapsto\pi$ injective on the convex slope ball: integrate
the quadratic form of its derivative along the segment between
two slopes.  The resulting open set $A_l$ of canonical initial
states $(z,\pi)$ on $s=s_0$ therefore satisfies
\begin{equation}\label{ex:initial-volume}
                c_1l^3\le |A_l|\le C_1l^3.
\end{equation}
Every initial and pre-stop momentum lies in a fixed bounded
set, since both $\widetilde g$ and $U$ are bounded on the
closed wider slope cone.

\emph{The stopped smooth flow and its exact Jacobian.}
From each state $x\in A_l$ follow the maximal geodesic in
the full $M$, using $s$ as parameter as long as it is in this
chart and $|\eta|<2\varepsilon_0$.  Stop also at $s=s_1$ or
at a box face.  Before these stops,
\[
            |z(s)-z_c|\le(\rho_0+4\varepsilon_0)l<l,
\]
so no spatial box face can be reached; the time interval is
also strictly inside the enlarged box.  For $s<s_1$ let
$D_s\subset A_l$ consist of the initial states whose arcs
exist through $s$ with all the pre-stop inequalities strict.
Smooth dependence in the interior of $M$ implies that $D_s$
and the joint pre-stop domain in $(s_0,s_1)\times A_l$ are open.

The hitting map $\Phi_s:D_s\to T^*\{s=\mathrm{constant}\}$,
in coordinates $(z,\pi)$, is smooth and injective.  Indeed
the geodesic equations are smooth on $M$, $\dot s>0$, and
backward uniqueness recovers the same initial state from a
hit before any stop.  Its canonical Jacobian is exactly one.
For completeness, put
$\mathcal H=\tfrac12g^{\alpha\beta}\pi_\alpha\pi_\beta$.
On the future unit shell $\mathcal H=-\tfrac12$,
$\partial\mathcal H/\partial\pi_0=\dot s>0$; hence
$\pi_0=-h(s,z,\pi)$ is a smooth root.  Implicit
differentiation of the shell equation gives the reduced
Hamilton equations
\[
        \frac{\dd z^i}{\dd s}=\frac{\partial h}{\partial\pi_i},
        \qquad
        \frac{\dd\pi_i}{\dd s}=-\frac{\partial h}{\partial z^i}.
\]
Their divergence in $(z,\pi)$ is zero, so
\begin{equation}\label{ex:section-volume}
        \Phi_s^*(\dd^3z\,\dd^3\pi)
                  =\dd^3z\,\dd^3\pi.
\end{equation}
Equivalently, the restriction of the canonical symplectic
form to the unit shell has the geodesic-flow direction as
kernel; the additional derivatives of a variable hitting
time disappear from its pullback.  This proves the same
volume statement for every transverse section used below.

Let
\[
       I(x)=\int_{\text{before stop}}
                    |\Gamma_M(\gamma_x(s))|\,\dd s.
\]
The joint open domain and the smooth interior flow make this
a measurable nonnegative function, possibly infinite.
Using \Cref{ex:section-volume}, Tonelli's theorem, and the
fixed bound on the momentum fibers gives
\begin{align}
 \int_{A_l} I(x)\,\dd x
 &=\int_{s_0}^{s_1}\int_{\Phi_s(D_s)}
                |\Gamma_M(s,z)|\,\dd^3z\,\dd^3\pi\,\dd s
       \notag\\
 &\le C\int_{Q_l(z_c)}|\widetilde\Gamma|\,\dd s\,\dd z
       =o(l^3).\label{ex:flow-average}
\end{align}
The equality of the weak and smooth Christoffel arrays is
only an almost-everywhere equality on $M$.  This suffices:
the same phase change of variables and Fubini's theorem show
that its exceptional chart-null set has zero time measure
on almost every initial arc.  Thus \Cref{ex:flow-average}
does not prescribe a trace of a Sobolev function on a curve.

\emph{Small integrals prevent the slope stop.}
Put $V=(1,\eta)$.  The smooth geodesic equation in $s$ time
is
\begin{equation}\label{ex:slope-equation}
 \frac{\dd\eta^i}{\dd s}
       =-\Gamma^i_{\alpha\beta}V^\alpha V^\beta
         +\eta^i\Gamma^0_{\alpha\beta}V^\alpha V^\beta.
\end{equation}
Consequently $|\dd\eta/\dd s|\le C_2|\Gamma_M|$ before
the stop.  By \Cref{ex:initial-volume,ex:flow-average}, for
small enough $l$ a set $G_l\subset A_l$ of measure at least
$|A_l|/2$ satisfies
$I(x)<\varepsilon_0/(2C_2)$.  On these arcs the slope remains
less than $3\varepsilon_0/2$.  It cannot reach the wider
threshold $2\varepsilon_0$.  The position estimate already
excludes a chart or box-face stop.

\emph{A surviving arc would give a forbidden bump.}
Suppose such an arc exists inside $M$ through $s_1$.  Choose
$\beta(s)$ supported in $(s_0,s_1)$, with
$\beta(s_c)=1$ and $|\beta'|\le C_\beta/l$, and put
\[
 \Delta z=z_c-z(s_c),\qquad
 z_r(s)=z(s)+r\beta(s)\Delta z,\qquad 0\le r\le1.
\]
Here $|\Delta z|\le(\rho_0+2\varepsilon_0)l$.
By \Cref{ex:cone-choice}, every graph $(s,z_r(s))$ has
spatial slope less than $c$.  The graphs stay in the ambient
rectangle and have the same two endpoints in $M$.  Their
first member is the assumed interior geodesic arc, whereas
the last passes through $(s_c,z_c)\in\partial M$.
This contradicts \Cref{ex:homotopy}.

It follows that every arc in $G_l$ has a maximal pre-stop
height $s_*\le s_1$, with no slope or chart stop.  Unit
normalization gives
\begin{equation}\label{ex:proper-time}
            0<c_3\le\frac{\dd s}{\dd\tau}=N\le C_3.
\end{equation}
The elapsed proper time to $s_*$ is therefore finite.  The
bounded spatial slope gives an ambient position limit in a
fixed compact subrectangle of the chart.  If this limit were
in $M$, smoothness of the metric there, the bounded unit
velocity, and the geodesic equation would give a limiting
velocity and smooth continuation.  If $s_*<s_1$, that
contradicts maximality of the pre-stop interval.  If
$s_*=s_1$, it gives the just-excluded surviving arc.
Thus the limit lies in $\partial M$.  Any continuation in
the full $M$ would, under the continuous embedding, take
this boundary value at the joining proper time, which is
impossible.  The finite lifetime just proved is therefore
the full future lifetime.  Nothing here solves an ambient
geodesic equation at the boundary.

\emph{Returning to $\Sigma$ and controlling the earlier arcs.}
We spell out why all the starting states have a past
intersection with $\Sigma$.  A maximal timelike geodesic in
a smooth spacetime has no interior causal endpoint.  Indeed,
in a relatively compact smooth temporal chart about a
hypothetical endpoint, all causal slopes $\eta=\dd z/\dd s$
are bounded.  The time component of the geodesic equation
gives
\[
  \frac{\dd}{\dd s}\log\frac{\dd s}{\dd\tau}
       =-\Gamma^0_{\alpha\beta}(1,\eta)^\alpha(1,\eta)^\beta.
\]
Its right side is bounded and $s$ has a finite endpoint.
Thus $\dd s/\dd\tau$ has positive upper and lower bounds.
The affine duration to the endpoint is finite, the velocity
is bounded, and the smooth geodesic equation gives a velocity
limit and continuation.  The same argument with reversed
orientation handles a past endpoint.  Hence maximal timelike
geodesics are causally inextendible; the Cauchy property
forces them to meet $\Sigma$ exactly once.  The starts here
are in $I_M^+(\Sigma)$, so the intersection is in their past.

For each initial state in $A_l$, this past hit occurs in
finite proper time and is transverse.  Smooth dependence and
the implicit function theorem give an open neighborhood on
which the hit and its time are smooth.  Shrink this to a
precompact neighborhood on which the transverse-section hit
map is a diffeomorphism onto its image and the complete
earlier flow tube has compact closure in the interior unit
shell.  Local injectivity is sufficient: the original local
time slice need not have a globally single-valued intersection
with every geodesic.  Second countability supplies countably
many such smaller neighborhoods covering $A_l$.  One of them
meets $G_l$ in positive measure.  Restrict to that intersection.

The backward hit preserves canonical section volume by the
unit-shell argument following \Cref{ex:section-volume}.
Identifying the future unit shell over $\Sigma$ with
$T\Sigma$ by the velocity
$\sqrt{1+h(w,w)}\,n+w$, its tangential covector is exactly
$h(w,\cdot)$.  The image $S$ is measurable and has positive
section measure.  The compact earlier flow tube is covered
by finitely many ambient charts with closures in $M$.
Taking smaller coordinate neighborhoods first gives a common
positive coordinate margin.  Compactness in the unit shell
then bounds the velocity components in the selected charts,
and bounds the earlier proper durations.  Each chart can be
taken precompact in a larger chart, so its metric, inverse,
and Christoffel $L^2$ norm are finite.  Add the single exit
chart, where \Cref{ex:proper-time} and the narrow cone supply
the velocity bound and the compact subrectangle supplies the
margin.  Taking a maximum of the bounds and a minimum of the
positive margins proves \Cref{ex:chart-witness} from proper
time zero onward.
\end{proof}

The two fixed-cone deformations are illustrated in \Cref{ex:figure}.

\begin{figure}[htbp]
 \centering
 \begin{tikzpicture}[x=1cm,y=1cm,>=Stealth,
  every node/.style={font=\small},
  interior/.style={blue!65!black,thick},
  deformation/.style={orange!85!black,thick},
  boundary/.style={black!65,densely dashed,thick}]
 \begin{scope}
  \node[font=\small\bfseries] at (0,4.55) {A common bottom fixes the cone};
  \draw[->,black!45] (-1.7,-.15)--(-1.7,3.75) node[above] {$s$};
  \draw[->,black!45] (-1.7,-.15)--(1.7,-.15) node[right] {$z$};
  \draw[black!25] (-1.4,.2)--(1.3,.2);
  \node[left,font=\scriptsize] at (-1.7,.2) {$-d_0$};
  \draw[boundary] (-1.3,1.65)--(-.15,1.5)--(.55,3.4)--(1.3,3.5);
  \node[align=left,font=\scriptsize,anchor=west] at (.7,3.85)
       {hypothetical\\higher exit};
  \draw[interior] (.55,.2)--(.55,2.95);
  \draw[deformation] (.55,.2)--(-.15,1.5)--(.55,2.95);
  \fill (.55,.2) circle (1.5pt) node[right] {$p$};
  \fill (.55,2.95) circle (1.5pt) node[right] {$q$};
  \fill[orange!85!black] (-.15,1.5) circle (2pt);
  \node[anchor=east,align=right,font=\scriptsize] at (-.27,1.95)
       {lower exit\\$(f(z),z)$};
  \fill (.55,3.4) circle (1.5pt);
  \draw[black!35,dotted] (.55,.2)--(-.35,1.1);
  \draw[black!35,dotted] (.55,.2)--(1.45,1.1);
  \node[font=\scriptsize,anchor=west] at (.73,.72) {fixed cone};
  \node[font=\scriptsize,align=center] at (0,-.65)
       {The vertical reference is interior.\\The deformed path meets the first exit.};
 \end{scope}
 \begin{scope}[xshift=7cm]
  \node[font=\small\bfseries] at (0,4.55) {A small bump keeps endpoints fixed};
  \draw[->,black!45] (-1.7,-.15)--(-1.7,3.75) node[above] {$s$};
  \draw[->,black!45] (-1.7,-.15)--(1.7,-.15) node[right] {$z$};
  \draw[black!25,rounded corners=2pt] (-1.3,.15) rectangle (1.25,3.55);
  \draw[black!25,dotted] (-1.3,.5)--(1.25,.5);
  \draw[black!25,dotted] (-1.3,3.2)--(1.25,3.2);
  \node[left,font=\scriptsize] at (-1.7,.5) {$s_c-l$};
  \node[left,font=\scriptsize] at (-1.7,1.85) {$s_c$};
  \node[left,font=\scriptsize] at (-1.7,3.2) {$s_c+l$};
  \draw[boundary] (-1.3,2.05)--(-.6,1.95)--(0,1.85)--(.28,3.4)--(1.25,3.4);
  \draw[interior] (.4,.5) .. controls (.45,1.25) and (.65,2.4) .. (.57,3.2);
  \draw[deformation] (.4,.5) .. controls (.43,.95) and (-.13,1.4) .. (0,1.85)
       .. controls (.1,2.25) and (.64,2.65) .. (.57,3.2);
  \fill (.4,.5) circle (1.5pt);
  \fill (.57,3.2) circle (1.5pt);
  \fill[orange!85!black] (0,1.85) circle (2pt);
  \draw[->,orange!85!black,thin] (.53,1.85)--(.06,1.85);
  \node[font=\scriptsize,anchor=south] at (.31,1.87) {$\Delta z$};
  \node[font=\scriptsize,anchor=west,blue!65!black] at (.7,2.55) {assumed surviving arc};
  \node[font=\scriptsize,align=center] at (0,-.65)
       {$|\Delta z|=O((\rho_0+\varepsilon_0)l)$;\\the slope change stays inside the cone.};
 \end{scope}
\end{tikzpicture}

 \caption{Two uses of the fixed timelike cone, with one spatial
 coordinate shown.  Left: a common bottom and a point below the
 higher first exit would admit a path through the lower exit if
 the graph were too steep.  Right: a sufficiently narrow
 geodesic surviving across a small graph box can be deformed,
 with its endpoints fixed, through the central exit.  Both
 deformations contradict the intrinsic-diamond argument; the
 deformed paths use only the fixed timelike cone in the ambient chart.}
 \label{ex:figure}
\end{figure}

\subsection{The existential tests and their countable cover}

We have obtained uniform chart witnesses along a positive-measure
family of entire future geodesics. The remaining step is to localize
that family in one small branch tube, where a finite packet can affect
all observations. This localization must retain enough phase volume
for the two data sets compared later.

Return to the near-Kerr developments.  Use the single
$p_{10}$ neighborhood fixed in \Cref{prop:bg-package}, and
write $\mu_d$ for the section measure on $T\Sigma$ just
defined.  Fix a countable basis $\mathcal B$ of coordinate
balls in $T\Sigma$, with compact closures in members of a
countable atlas.  The centers and radii can be taken rational
in these coordinates, so the basis is fine enough for
Lebesgue differentiation with any smooth positive density.

\begin{definition}[Necessary existential extension tests]
\label{def:exit-test}
For $O\in\mathcal B$ and $B\in\N$, $B\ge1$, let
$E(O,B)\subset\Ueps$ consist of the data satisfying both
of the following requirements.
\begin{enumerate}[label=(\roman*)]
\item There is an extension of $M_d$ of the regularity and
      type specified in Definition~\ref{intro:extension}, and a measurable
      $S\subset O$ with
      \begin{equation}\label{ex:density-test}
                       \mu_d(S)>.99\,\mu_d(O),
      \end{equation}
      such that every seed in $S$ has finite full future
      lifetime.  There are at most $B$ precompact ambient
      charts $\psi_j:U_j\to V_j$ witnessing
      \Cref{ex:chart-witness} along each entire such geodesic,
      from $\Sigma$ to its full endpoint.  On the whole of
      every $V_j$, each of the metric norm, inverse metric
      norm, and Christoffel $L^2$ norm is at most $B$.
      At each point a selected chart has velocity norm at
      most $B$ and Euclidean distance to
      $\R^4\setminus V_j$ at least $1/B$.
\item There is an observation neighborhood $\mathcal N$ satisfying all three
      parts of \Cref{lem:bg-branch-tube} for this datum, with
      $\overline O\subset\mathcal N$ and the optical ordering chosen so that
      $v\to\infty$.
      One choice of its central branch track, compact state ranges,
      positive $p$ bounds, transverse earlier flow tube, interval $J$,
      observation boxes, admissible launch patch, unbounded sequence of
      $T$, and constants $u_2,c_2,c_0$ serves every seed in $\overline{\mathcal N}$.
      In particular,
      the entry-to-observation portions fit
      \eqref{bg:experiment-domain} with its outer margins and positive
      slice-length bound. The admissible patch has the low-order ray
      and collar properties specified before \Cref{lem:bg-branch-tube}.
\end{enumerate}
All tube constants, charts, patches, and subsequences in this
definition are existential witnesses for the given datum.
They are not additional countable indices and are not
required to work in a neighborhood of that datum.
\end{definition}

The measurable good family in this definition may depend on
the extension.  For later averaging, whenever its point and
velocity satisfy several chart conditions, select the least
chart index.  The finite collection of coordinate margin and
velocity conditions is Borel on the interior unit shell, so
this gives a measurable selector, also when it depends on the
momentum.  No closedness assertion about $E(O,B)$ is made.

\begin{proposition}[Countable cover of all future extensions]
\label{prop:exit-cover}
Every datum in $\Ueps$ admitting a future
$C^0\cap W^{1,2}_{\mathrm{loc}}$ extension belongs to some
$E(O,B)$.  Consequently the extendible set is precisely
\[
                  \bigcup_{O\in\mathcal B}\ 
                         \bigcup_{B=1}^{\infty}E(O,B).
\]
This cover includes extensions whose actual exit occurs
after first runout from the comparison development.
\end{proposition}

\begin{proof}
Fix such a datum and extension.  \Cref{prop:exit-family}
provides a measurable family $S_*$ of positive $\mu_d$ measure,
finite full lifetimes, and one finite set of ambient chart
witnesses.  By \Cref{prop:bg-runout}, every seed in $S_*$
crosses the comparison cylinder and its first comparison
runout is a branch or a corner.  The corner subset is null.
The inputs here are precisely the finite-full-lifetime
classification of \cite[\CompanionRef{foundation}{cmp:classification}]{KerrCTwo} and the
section-volume corner removal of
\cite[\CompanionRef{signed}{test:corner-null}]{KerrCOne}; the latter uses only the
low-order corner estimates of
\cite[\CompanionRef{foundation}{cmp:corner-ratios}]{KerrCTwo}.  These statements have neither a
regular-extension hypothesis nor a curvature bound along
the selected seed.

Choose a density point $w$ of $S_*$ outside that null subset,
with $w\in S_*$.  It is a branch seed.  Apply
\Cref{lem:bg-branch-tube} and, if necessary, interchange the
optical labels.  This gives an open seed neighborhood $\mathcal N$ of
$w$ with all the tube, footprint, and finite-domain
properties in \Cref{def:exit-test}(ii). These are consequences of
the background estimates and hold independently of the extension
chart witnesses.

By density and the fine countable basis, choose $O\ni w$
with $\overline O\subset\mathcal N$ and
$\mu_d(S_*\cap O)>.99\mu_d(O)$.  To justify using countably
many balls, small balls centered at a density point can be
enlarged by an arbitrarily small relative amount to balls
with rational center and radius.  Smooth positivity of the
density keeps their measures comparable, so their relative
complement measure still tends to zero.  Set $S=S_*\cap O$.
Choose an integer $B$ at least the number of the ambient
charts, all their bounds, and the reciprocal of their common
margin.  This proves $d\in E(O,B)$.

Every observation invoked here lies strictly before first
runout from the comparison region.  The finite full lifetime
and ambient chart witnesses were obtained independently for
the entire geodesic in $M_d$.  Thus first runout is used only
to locate earlier observations; it need not be the ambient
exit.  Finally the reverse inclusion in the displayed union
is part of \Cref{def:exit-test}(i).
\end{proof}
